% Working title — alternates in papers/methodology/abstract-and-claims.md
% "Contracts at Machine Speed: An Inverted SDLC for an Agent-Operated Software Factory"
% Genre: experience report. Preprint-first (arXiv cs.SE/cs.AI/cs.DC + Zenodo) -> ICSE SEIP / FSE Industry.
% Build: pdflatex main.tex && pdflatex main.tex   (or `make`, or Tectonic, or Overleaf)
% Preamble copied verbatim from ../../arxiv/main.tex (the protocol paper) — keep in sync.
\documentclass[11pt]{article}
\usepackage[margin=1in]{geometry}
\usepackage[T1]{fontenc}
\usepackage{lmodern}
\usepackage{microtype}
\usepackage{amsmath,amssymb}
\usepackage{graphicx}
\usepackage{tikz}
\usetikzlibrary{arrows.meta,positioning}
\usepackage{float} % [H] pins tables/figures inline with their sections
\usepackage{array}
\usepackage{enumitem}
\usepackage{listings}
\usepackage[hidelinks]{hyperref}
\usepackage{url}
\urlstyle{same}
% Figures compile with OR without the converted PDFs (convert figures/methodology/*.svg).
\newcommand{\paperfig}[2]{\IfFileExists{#1}{\includegraphics[width=\linewidth]{#1}}%
{\fbox{\parbox{0.95\linewidth}{\centering\vspace{1.4em}\itshape #2\\[2pt]\textnormal{[figure placeholder --- draw as TikZ or convert \texttt{figures/methodology/*.svg} to PDF]}\vspace{1.4em}}}}}
\newcommand{\TODO}[1]{\textbf{\textcolor{red}{[TODO: #1]}}\typeout{TODO: #1}}
\usepackage{xcolor}
\lstset{basicstyle=\ttfamily\small,frame=single,breaklines=true,columns=fullflexible}

\title{Contracts at Machine Speed:\\ An Inverted SDLC for an Agent-Operated Software Factory}
\author{David S. Tufts\\
MyndHyve Inc.\\
\emph{\small Project steward (see Positionality, Section~1)}\\
Grand Rapids, Michigan, USA\\
\url{https://davidtufts.me}\quad\url{https://www.linkedin.com/in/davidtufts}}
\date{2026}

\begin{document}
\maketitle

\begin{abstract}
Autonomous AI agents now perform much of the work in a software life cycle, but the
literature on agentic software engineering treats agent coordination as a problem \emph{within}
a single development workflow. The unsolved problem is coordination \emph{across} independent
services: when an agent building one feature needs another service to change its interface.
We report on a methodology in which autonomous agents propose, negotiate, and adopt
versioned wire-contract changes across a multi-service estate---to our knowledge the first
end-to-end account of agent-driven contract evolution at machine speed. Each shared service
publishes a versioned contract; embedded per-service architect agents negotiate changes over
a coordination bus through a governed RFC process (Draft~$\rightarrow$~Active~$\rightarrow$~Accepted,
with risk-scaled review windows); and every change is verified by a conformance suite. A second
mechanism makes the verification trustworthy---\emph{honesty by construction}: a service advertises
a capability only when it is reachable in live state, and a strict conformance mode mechanically
fails any advertised-but-unimplemented claim. The work sits within an \emph{inverted SDLC} in which
agent compression of Create and Operate moves the binding effort to Plan and Validate, and an
operating model of mission-focused Builder Teams whose shared functions act as Contract Guardians by
embedding standards into the contracts themselves. We present a retrospective over a governance
corpus of 110 RFCs and 164 architecture decision records,
demonstrate the honesty-enforcement mechanism reproducibly, and trace contract changes graduating
across two distinct hosts on a dual-witness conformance bar, releasing the corpus and harness
templates as an artifact. As a single-steward experience report --- the two hosts share change
control --- we make no independent-validation, causal, or efficacy claim, and we present economics
only as a transparent, non-causal model.
\end{abstract}

\section*{1. Introduction}
In a conventional software life cycle most engineering effort is spent operating and maintaining what
already exists; creating new code is a minority of the work, and planning and validation a smaller
minority still. Large-language-model agents invert this distribution. As agents absorb both the
writing of code and much of the operational toil, the binding constraint migrates to the front of the
life cycle --- to deciding \emph{what} to build and \emph{verifying} that what was built conforms to
its specification. We call this the \emph{inverted SDLC}. The framing itself is not new: the Shift-Up
framework argues that machine-readable requirements, architecture models, and architecture decision
records act as guardrails that stabilize agent behavior and shift human effort toward design and
validation [SHIFTUP], and surveys of ``vibe coding'' describe developers validating agent output by
observing outcomes rather than reading code line by line [VIBE]. Our concern is what makes that shift
\emph{safe across more than one service}.

The agentic-software-engineering literature is substantial but treats agent coordination as a problem
\emph{within} a single development workflow: planning, tool use, and multi-step execution inside one
repository or task [TOSEM, AGENTPROG]. The problem we address is orthogonal and, to our knowledge,
unaddressed: coordination \emph{across} independent services, when an agent building one feature needs
another service to change its published interface. Current agent-interaction protocols do not fill
this gap --- across MCP, A2A, and related schemes agents discover capabilities and select a protocol
version, but none provide a mechanism to propose, negotiate, or version \emph{changes} to a shared
service contract [A2A, AGENTPROTO]. We report on a methodology in which autonomous agents propose,
negotiate, and adopt versioned wire-contract changes across a multi-service estate, governed by a
Request-for-Comments (RFC) process the agents themselves execute, and verified by a conformance suite
whose strict mode makes a dishonest capability claim mechanically fail. We describe the architecture,
the operating model, and the agentic harness; we present a retrospective over a real governance
corpus of 110 RFCs and 164 architecture decision records; we demonstrate the honesty mechanism
reproducibly; and we trace contract changes graduating across two distinct hosts.

\paragraph{Terminology.} Three terms in this area are overloaded, and we fix our usage now. By
\emph{contract} we mean a versioned service or API \emph{wire contract} --- the externally observable
interface between services --- not the resource-governance ``agent contract'' (token, time, and budget
bounds on execution) of the multi-agent-systems literature [AGENTCONTRACT]. By \emph{AI software
factory} we mean a \emph{methodology} for autonomous agents building features across a multi-service
estate; this is distinct from the vendor sense of ``AI factory'' as model-serving infrastructure
[NVIDIA] and from the DoD / Platform One ``software factory'' as a DevSecOps pipeline. The nearest
published neighbor is the ``agentic SDLC'' / ``agent factories'' framing [MCKINSEY], which denotes
central capabilities that \emph{produce} agents rather than an estate the agents \emph{evolve}.

\paragraph{Positionality.} This manuscript is authored by the steward of the estate under study, not
an independent third party. The protocol corpus, the conformance suite, and both host implementations
are steward-produced, and a substantial fraction of the corpus --- and of this manuscript --- was
produced by the very agents the paper studies, under human review. We therefore treat all material as
evidence of \emph{implementability and internal consistency}, never as third-party validation or
adoption, and we flag throughout where a claim would require an independent party to establish. We
disclose the use of AI assistance explicitly because it is also our subject; Section~14 states which
artifacts were agent-produced and how they were reviewed. This stance is a deliberate
threat-to-validity control, stated up front so the evidence below is read correctly.

\section*{2. Problem: Coordinating Change in an Agent-Built Estate}
When integration knowledge lives in people's heads and in code review, a service's consumers learn of
an interface change through conversation, convention, and the occasional broken build. In an estate
where agents propose and apply changes faster than humans can review them, that implicit knowledge
fails in three ways, each acute at machine speed:
\begin{itemize}[leftmargin=*]
  \item \textbf{Silent contract drift.} A service changes its wire shape and its consumers discover
        the break in production, because nothing forced the change to be declared against a shared,
        versioned contract.
  \item \textbf{Dishonest capability claims.} A service advertises a capability it does not actually
        honor --- through optimism, staleness, or a half-finished change --- and integrators build
        against a fiction.
  \item \textbf{Uncoordinated rollouts.} A breaking change ships before consumers can adopt it, or
        consumers are never told it exists.
\end{itemize}
These are survivable when a human reviewer sits in the path of every change; they are not survivable
when the proposer is an agent and the reviewer would have to be one too. The methodology in this paper
is organized around a single response: \emph{make the contract the locus of coordination, and make
conformance to it continuous and automatic} --- so that ``validate'' is no longer ``did the output
look right?'' but ``does this implementation still satisfy every contract it touches?'' The remaining
sections develop that response and the evidence for it.

\section*{3. Contribution}
% The C1--C10 claims (see abstract-and-claims.md), ORDERED BY NOVELTY (lead with C1/C2).
This paper makes the following contributions, each tagged with its evidence tier
(\textbf{D}emonstrated / \textbf{O}bserved / \textbf{A}rgued):
\begin{enumerate}[leftmargin=*]
  \item \textbf{(C1, novel)} Agent-negotiated, cross-service \emph{versioned-contract evolution}:
        embedded per-service architect agents propose and negotiate changes to versioned wire
        contracts over a coordination bus, under a governed RFC lifecycle. We are not aware of prior
        work in which agents negotiate versioned \emph{service-contract} changes; the agent-protocol
        literature provides capability discovery and version selection but no contract-change
        negotiation [A2A, AGENTPROTO] (Section~4).
  \item \textbf{(C2, novel)} An RFC change process (comment windows, status lifecycle) executed \emph{by}
        agents at machine speed.
  \item \textbf{(C3, demonstrated --- not a ``first'' claim)} \emph{Honesty by construction}:
        advertisement bound to live state + strict conformance fails dishonest claims. The
        fail-on-dishonest \emph{gate} is anticipated by model-based publication gates [RELWS] and
        bi-directional contract testing [PACTFLOW]; our increment is binding advertisement to
        \emph{continuous live runtime reachability} within an autonomous multi-agent
        contract-evolution estate, demonstrated reproducibly (Section~10).
  \item \textbf{(C4--C9)} The contract gate (``friction partner''); the ADR-local/RFC-external split
        (cite + out-distance [SHIFTUP]); the reproducible governance corpus; the Builder-Team /
        Contract-Guardian operating model; the agentic harness; the inverted-SDLC framing.
  \item \textbf{(C10, model)} A transparent, non-causal economic framing (Section~12).
\end{enumerate}
Claims C1, C2, C4, and C6 are evidenced by a released, re-runnable governance corpus
(Sections~10--11); C3 is demonstrated by a reproducible experiment (Section~10); C5 and C7--C9 are
argued from the design and the operating record. We make no controlled-efficacy, generalization, or
causal-ROI claim: the estate has a single change-control authority (Section~14), and these are
deferred to a validation agenda (Section~15).

\section*{4. Related Work}
% Anchored by related-work-and-novelty.md. Flag source tiers (peer-reviewed vs preprint vs vendor).
\TODO{Four bodies: (A) agentic SE [TOSEM, AGENTPROG, MICROSVC, APIFIRST] --- coordination is
intra-workflow; (B) contract/API governance [TYPESAFE, CDVO] --- largely AI-free baselines;
(C) the inverted-SDLC / spec-driven frame [SHIFTUP, VIBE]; (D) agent infrastructure
[IOA, AGENTIC-SC, AGENTCONTRACT]. Carve the delta against [SHIFTUP] explicitly (it reframes ADRs as
guardrails but has no local/external split, no wire contracts).}

\section*{5. The Operating Model: Builder Teams in an Architect Mesh}
\TODO{Builder Teams (Product Owner / UX Designer / AI Engineer); shared functions as Contract
Guardians embedding standards as contract validation rules, not review queues. Team-topologies /
platform-engineering lineage.}

\section*{6. The Process: Discovery and the Build Loop}
\TODO{Discovery (deep research -> multi-persona PRD lenses {Spec, Schema, Security, Conformance,
Compatibility} -> contract-as-friction -> dependency-aware roadmap -> ADR/RFC split -> architect
ratification). Then the Build Loop as a listing:}
\begin{lstlisting}[caption={The build loop (per ADR + ratified RFC, in dependency order).},label={lst:loop}]
for each ADR and ratified RFC (in dependency order):
  for each implementation phase:
    architect       # check phase design vs built code and local ADRs
    goal            # autonomous implementation of a feature or contract clause
    verify-contract # automated check of implementation vs contract schema
    ux-review       # check UI vs design-system contract
    merge           # merge behind a feature toggle
\end{lstlisting}

\section*{7. The Architecture}
\TODO{Local feature lifecycle (roadmap -> ADR -> feature package -> catalog/toggle) ->
the contract gate (three buckets; the ``friction partner'') -> contract governance (RFC lifecycle,
risk-scaled windows) -> the architect mesh + coordination bus.}
\begin{figure}[H]\centering
  \paperfig{figures/methodology/fig-contract-gate.pdf}{The contract gate: host-extension / rides-accepted / touches-the-wire, wired to a failing test.}
  \caption{The contract gate as a forcing function.}\label{fig:gate}
\end{figure}

\section*{8. Honesty by Construction}
\TODO{Validate \emph{as} continuous contract conformance: the discovery $\leftrightarrow$ live-state
$\leftrightarrow$ strict-conformance loop. This section hosts Study~B (Section~10).}

\section*{9. The Agentic Harness}
\TODO{Context-as-contracts, skills-as-consumers, hooks-as-obligations, memory. The transferable
artifact; ship templates with the corpus release. Cite AGENTS.md/CLAUDE.md-style context conventions
and [SHIFTUP] as lineage.}

\section*{10. Evidence: Corpus Retrospective and the Honesty Experiment}
\paragraph{Study A (corpus retrospective; \emph{Observed}).} A deterministic analysis of the
governance corpus (script + per-item data in \texttt{evidence/corpus-analysis/}; snapshot
\texttt{openwop@d7a635b6}, \texttt{openwop-app@bfcddc7d}). The RFC corpus is $n=110$ (107 Accepted,
2 Active, 1 Draft), of which 104 record at least one lifecycle transition. The decision-record corpus
is $n=164$ (144 implemented, 11 Accepted, 6 Superseded, 2 Proposed, 1 Withdrawn); 64 (39\%) carry a
formal correction note, quantifying the ``correct, don't rewrite history'' discipline. The headline
result is the \emph{contract-gate distribution}. A first-pass whole-text heuristic over-counted
wire-touching changes; a validation pass that classifies each record from its \emph{explicit}
RFC-gate declaration --- with the residual hand-coded (script and audit in
\texttt{evidence/corpus-analysis/gate-audit.md}) --- corrects it: of 159 classifiable decision
records, \textbf{155 (97.5\%) required no new wire RFC} (120 host-extension, 35 riding an
already-accepted RFC) versus only \textbf{4 (2.5\%) wire-touching}. (The heuristic's
$\sim$8$\times$ over-count of the wire-touching bucket --- it matched ``new RFC'' inside ``\emph{no}
new RFC'' --- and the resulting 57\% heuristic/audit agreement are reported in full; the corrected
figure makes the gate \emph{more} selective than the heuristic suggested, evidencing C1/C2.)
\TODO{Add conformance growth (scenarios per accepted RFC) via git-history mining (point-in-time:
v1.37.0, 370 scenarios, 82 fixtures).}

\paragraph{Study B (honesty experiment; \emph{Demonstrated}).} A reference host whose
\texttt{workflowChainPacks} advertisement is recomputed from live execution state
(\texttt{existsSync} per \texttt{.well-known/openwop} request) is driven through three honesty
postures by where its pack registry points --- with \emph{no host-code modification} --- and the
gated behavioral conformance scenario is run in default and strict
(\texttt{OPENWOP\_REQUIRE\_BEHAVIOR=true}) mode (harness + captured run in
\texttt{evidence/honesty-experiment/}):
\begin{center}\small
\begin{tabular}{llll}
\textbf{Posture} & \textbf{Advertises} & \textbf{Mode} & \textbf{Result}\\\hline
honest + implemented (backed) & yes & strict & \textbf{6/6 pass} (non-vacuous)\\
honest minimal (no backing) & no & default & 6/6 pass (honest skip)\\
claims coverage, advertises nothing & no & strict & \textbf{6/6 fail}\\
\textbf{dishonest} (empty registry) & \textbf{yes} & strict & \textbf{3/6 fail}\\
\end{tabular}
\end{center}
The decisive contrast is the last two honest rows versus the dishonest row: the host advertises
\texttt{workflowChainPacks:\{supported:true\}} but cannot expand (the behavioral leg returns
\texttt{pack\_not\_found}), so strict conformance \emph{fails}. \textbf{A host cannot advertise a
capability it does not deliver without failing its own conformance suite.} Per Section~4 this
fail-on-dishonest gate is anticipated by prior model-based publication gates and bi-directional
contract testing [RELWS, PACTFLOW]; the demonstrated increment is the binding of the advertisement
to \emph{continuous live runtime state}, re-evaluated per request, within the autonomous estate.

\section*{11. Evidence: Cross-Host Contract Evolution}
\paragraph{The estate (two hosts, one wire contract).} The methodology is exercised across two
genuinely distinct codebases and production deployments that advertise the same wire contract at
their \texttt{.well-known/openwop} discovery documents: the open reference host
\textbf{openwop-app} (\texttt{workflow-engine}, \texttt{app.openwop.dev}) and the steward's product
host \textbf{MyndHyve} (\texttt{workflow-runtime}, a separate Firebase codebase at
\texttt{api.myndhyve.ai}). A change graduates \texttt{Active}~$\rightarrow$~\texttt{Accepted} only
when a \emph{second} host advertises the capability and passes its gated conformance scenarios
\emph{non-vacuously} (\texttt{OPENWOP\_REQUIRE\_BEHAVIOR=true}) --- the \emph{dual-witness bar}.
Adoption is coordinated by separate per-repo agent sessions exchanging handoff artifacts over a
message bus; the agents implement, advertise, and run conformance on each host, and the steward is
the single change-control authority that flips status (\emph{agent-executed, steward-supervised}).
Full trace and citations in \texttt{evidence/cross-host-case/}.

\paragraph{Worked example --- RFC 0050 (SAML/SCIM).} Authored \texttt{Draft} 2026-05-24 from a real
product need; full wire surface landed and promoted to \texttt{Active} 2026-06-01; graduated
\texttt{Accepted} the same day when MyndHyve advertised
\texttt{auth.profiles:[openwop-auth-saml,openwop-auth-scim]} and passed the gated legs
\textbf{19/19 non-vacuously} (the count rose 13$\rightarrow$19 once the IdP engaged the behavioral
legs), with all seven SAML negative variants (incl.\ signature-wrapping) rejected live. At scale the
record is routine: one cohort graduated \textbf{8 RFCs Draft$\rightarrow$Accepted in a day} on a
single MyndHyve revision (28~PASS / 0~FAIL), and recent changes (RFC 0095, 0099, 0100, 0102, 0103,
0105/0106) graduated as dual-witness pairs.

\paragraph{Honest non-adoption (the gate is real).} The same estate records \emph{retractions} and
\emph{opt-outs}: MyndHyve advertised one capability \texttt{(maxNodeExecutions)} in error and
\emph{honestly retracted} it; several RFCs sit at documented host opt-outs; a fail-loud host correctly
\emph{rejected} malformed packs. A dishonest advertisement is removed, not left standing --- the
honesty-by-construction loop (Section~10) observed at estate scale.

\paragraph{Study D (velocity; \emph{Observed}).} Mining the dates recorded in each RFC's
lifecycle field (script + per-RFC data in \texttt{evidence/operational-stats/}), the entire
110-RFC program ran in \textbf{$\sim$54 days}, and over the 107 Accepted RFCs the authored-to-Accepted
span has a \textbf{median of 0 days} (mean 1.1, max 12): \textbf{56\% graduated same-day}, 76\% within
one day, 97\% within seven. This is the machine-speed signal --- but it must be read with two
conditions. First, it is enabled by \emph{single-maintainer governance with waived comment windows}
(Study A: 40 RFCs note a waived/bootstrap window); the risk-scaled windows (Section~8) are the designed
brake for the multi-reviewer case, which is not yet exercised, so the honest claim is that the
methodology removes the \emph{implementation/integration} bottleneck, not that governed review takes
zero time. Second, a span is a \emph{recorded-activity} window, not human-effort time, and we make no
human-hours-saved claim from it (Section~12). \TODO{Draw Fig.~\ref{fig:seq} as a TikZ swimlane.}
\begin{figure}[H]\centering
  \paperfig{figures/methodology/fig-change-sequence.pdf}{spec -> openwop-app -> MyndHyve: a wire change graduating on the dual-witness bar.}
  \caption{Cross-host contract evolution: a change graduates only on a second host's non-vacuous conformance.}\label{fig:seq}
\end{figure}

\section*{12. Economic Framing (Descriptive, Non-Causal)}
% QUARANTINED. Glass-box model, n, measured-vs-modeled, confounds inline. NOT in the abstract headline.
% RESOLVED (R3): the "11.6:1 / $86-hr / workday-a-week" figures are UNATTRIBUTABLE -> dropped.
% Present ONLY a transparent internal model bracketed by the citable independents.
\TODO{If an internal value-recovery model is shown, state its formula, n, and window in full and label
it non-causal; name confounds (selection bias, no control, blended-rate abstraction, the perception
gap [METR] demonstrated). Do NOT reproduce the unsourceable vendor figures. Bracket strictly with the
independent evidence: METR RCT [METR] ($-19\%$, and the proof self-reported productivity is unreliable)
and the peer-reviewed Management Science 3-RCT [MSRCT] ($+26.08\%$ tasks, 4{,}867 devs).}

\section*{13. Claims and Evidence Discipline}
\TODO{The demonstrated / observed / argued ladder; map every claim in Section~3 to its tier.}

\section*{14. Limitations and Threats to Validity}
\paragraph{Single change-control authority (the central threat).} The estate spans two genuinely
distinct codebases and deployments (Section~11), which is real cross-host portability evidence ---
but both are operated by the same steward. In the project's own evidence-tier taxonomy this is
\textbf{Tier~2 (steward-affiliated sibling host)}, explicitly \emph{not} Tier~3 (independent
organization); no Tier-3 host exists. We therefore present the dual-witness records as evidence of
cross-codebase \emph{implementability}, never as independent third-party validation. (Historical RFC
and changelog wording calls the sibling host ``non-steward''; the governance taxonomy later corrected
that label and we use the corrected one.) A Tier-3 host plus independent external review is the
decisive next evidence --- deferred to the Validation Agenda (Section~15).
\TODO{Also: AI-authored corpus + AI-assisted manuscript (disclose); self-evaluation and selection
bias; the perception-reality gap demonstrated by [METR]; the contract-gate metric is a prose-mined
heuristic pending a manual coding pass; the honesty-gate mechanism is anticipated by [RELWS, PACTFLOW]
and our contribution is its live-state-bound, multi-agent form; the economic model is non-causal with
no control and no claimed productivity figure is independently sourced.}

\section*{15. Validation Agenda}
\TODO{What independent adoption would test; pre-registered hypotheses; the second-host / second-team
experiments that would move claims from observed/argued to demonstrated.}

\section*{16. Reproducibility and Artifact Availability}
\TODO{The released governance corpus (scrubbed RFCs/ADRs/conformance history + redacted bus transcripts),
analysis scripts that regenerate Study~A's numbers, and the agentic-harness templates. CC BY + Zenodo
artifact DOI. Note the scrub boundary (secrets, tenant identifiers, internal hostnames).}

\section*{17. Conclusion}
\TODO{The wager: govern at machine speed by moving human judgment into the contract, so review becomes
continuous and automatic rather than a queue.}

\begin{thebibliography}{99}
\small
\bibitem[TOSEM]{tosem} J.~He, C.~Treude, D.~Lo. ``LLM-Based Multi-Agent Systems for Software Engineering: Literature Review, Vision and the Road Ahead.'' ACM TOSEM, 2025. arXiv:2404.04834.
\bibitem[AGENTPROG]{agentprog} Y.~Wang et al. ``AI Agentic Programming: A Survey.'' arXiv:2508.11126, 2025. \emph{(preprint)}
\bibitem[MICROSVC]{microsvc} ``Specification Complexity in Microservice-Based Applications.'' arXiv:2508.20119, 2025. \emph{(preprint)}
\bibitem[APIFIRST]{apifirst} ``From Specification to Service: Accelerating API-First Development Using Multi-Agent Systems.'' arXiv:2510.19274, 2025. \emph{(preprint)}
\bibitem[TYPESAFE]{typesafe} J.~Costa~Seco et al. ``Robust Contract Evolution in a Type-Safe MicroServices Architecture.'' \textlangle Programming\textrangle{} 2020. arXiv:2002.06185.
\bibitem[CDVO]{cdvo} ``A Compatibility-Driven Version Orchestrator for Microservice Contract Evolution.'' MDPI Digital 5(3):27, 2025.
\bibitem[SHIFTUP]{shiftup} V.~Lipsanen, T.~Mikkonen et al. ``Shift-Up: A Framework for Software Engineering Guardrails in AI-native Software Development.'' VibeX 2026. arXiv:2604.20436. \emph{(preprint; initial findings)}
\bibitem[VIBE]{vibe} Y.~Ge et al. ``A Survey of Vibe Coding with LLMs.'' arXiv:2510.12399, 2025. \emph{(preprint)}
\bibitem[IOA]{ioa} ``Internet of Agents.'' IEEE TCCN, 2025. arXiv:2505.07176.
\bibitem[AGENTIC-SC]{agenticsc} ``Agentic Services Computing.'' arXiv:2509.24380, 2025. \emph{(preprint)}
\bibitem[AGENTCONTRACT]{agentcontract} ``Agent Contracts: A Formal Framework for Resource-Bounded Autonomous AI Systems.'' COINE/AAMAS 2026. arXiv:2601.08815.
\bibitem[A2A]{a2a} Agent2Agent (A2A) Protocol Specification, v1.0, 2026. \url{https://a2a-protocol.org/latest/specification/}. \emph{(primary spec)}
\bibitem[AGENTPROTO]{agentproto} ``A Survey of AI Agent Protocols (MCP/ACP/A2A/ANP).'' arXiv:2505.02279, 2025. \emph{(preprint)}
\bibitem[METR]{metr} METR. ``Measuring the Impact of Early-2025 AI on Experienced Open-Source Developer Productivity'' (RCT). 2025. arXiv:2507.09089. \emph{(independent)}
\bibitem[MSRCT]{msrct} Z.~Cui, M.~Demirer, S.~Jaffe, L.~Musolff, S.~Peng, T.~Salz. ``The Effects of Generative AI on High-Skilled Work: Evidence from Three Field Experiments with Software Developers.'' Management Science, INFORMS, 2026. DOI 10.1287/mnsc.2025.00535. \emph{(peer-reviewed; 3 RCTs, 4{,}867 devs; vendor tool)}
\bibitem[RELWS]{relws} E.~Ramollari, D.~Dranidis, A.~J.~H.~Simons. ``Reliable Web Service Publication and Discovery through Model-Based Testing and Verification.'' Univ.\ of Sheffield. \emph{(peer-reviewed; X-machine publication gate --- closest prior art for the honesty gate)}
\bibitem[PACTFLOW]{pactflow} SmartBear/PactFlow. ``Bi-Directional Contract Testing'' (compatibility checks; \texttt{can-i-deploy} deploy gate). \emph{(product docs)}
\bibitem[MCKINSEY]{mckinsey} McKinsey. ``Rewiring software delivery for the agentic era.'' 2025. \emph{(vendor/consultancy; ``agentic SDLC''/``agent factories'')}
\bibitem[NVIDIA]{nvidia} NVIDIA. ``AI Factory'' white paper (agentic AI control plane). \emph{(vendor; infrastructure sense)}
\end{thebibliography}

\end{document}
